\section*{ЗАКЛЮЧЕНИЕ}
\addcontentsline{toc}{section}{ЗАКЛЮЧЕНИЕ}

В данной выпускной квалификационной работе была проведена разработка интеллектуальной системы, обеспечивающей повышение эффективности производственного планирования посредством адаптивной диспетчеризации операций с использованием нейро-сетевых методов. Актуальность выбранной темы обусловлена необходимостью оперативного составления производственных расписаний в условиях гибкого дискретного производства, где классические MILP-решатели требуют значительных вычислительных ресурсов, а стандартные ERP/MES-системы не способны автоматически адаптироваться к текущему состоянию оборудования.

Основные результаты, полученные в ходе выполнения работы:
\begin{enumerate}
\item Разработан модуль управления заказами и производственными ресурсами, учитывающий технологические типы операций и связи «ресурс–тип». Реализован графический интерфейс с формами создания и редактирования заказов, ресурсов и их параметров, а также дашборд для мониторинга очередей и drag-and-drop управления задачами.
\item Разработана подсистема построения производственных расписаний, использующая эвристический алгоритм EDD + эффективная загрузка и нейросетевые диспетчеры PointerNet, SLIM, PPO. Обеспечена визуализация расписания в виде диаграммы Ганта и возможность выбора метода планирования.
\item Реализован модуль обучения нейросетевых моделей на синтетических данных, генерируемых по правилу минимальной эффективной загрузки. Модуль поддерживает три парадигмы обучения – поведенческое клонирование, самоконтролируемое обучение с саморазметкой и обучение с подкреплением, а также сохраняет веса, метрики и отчёты о качестве классификации.
\item Проведён сравнительный анализ четырёх методов планирования на стандартизированном тестовом сценарии (120 заказов, 5 ресурсов с ограничениями и различной готовностью). Установлено, что модель SLIM обеспечивает наименьшее среднее опоздание (8,1\,ч), минимальный makespan (35,0\,ч) и наилучшую равномерность загрузки (коэффициент вариации 0,171), превосходя чистую эвристику и другие нейросетевые методы.
\end{enumerate}

Практическая ценность работы заключается в том, что разработанная система позволяет сократить время построения расписаний, повысить равномерность загрузки оборудования и уменьшить опоздание заказов за счёт применения обученных нейросетевых диспетчеров, что в конечном счёте ведёт к снижению производственных издержек и штрафов за срыв сроков.

Научно-практическая новизна работы заключается в создании единой программной среды, объединяющей эвристические и нейросетевые методы планирования, а также в экспериментальном доказательстве преимущества саморазмеченного обучения (SLIM) над классическим поведенческим клонированием в условиях ограниченной совместимости ресурсов и операций.

Перспективами дальнейшего развития проекта являются расширение системы для поддержки динамических событий в реальном времени (поломки оборудования, срочные заказы), интеграция с существующими MES/ERP-системами.
